\section{MI-07: no finite two-unitary constant for the maximal symmetric modulus}
\label{sec:mi07}

For a square matrix $Z$, define its maximal symmetric modulus by
\[
 \maxmod(Z)=\lim_{r\to\infty}\bigl(|Z|^r+|Z^*|^r\bigr)^{1/r},
\]
where $r$ runs through positive integers. This is the spectral-order join $|Z|\vee|Z^*|$. The target is the constant-one assertion of \cite[Conjecture 3.12]{BL2026}.

\begin{theorem}\label{thm:mi07}
There is no finite nonnegative constant $C$, even in dimension two, such that every $A,B\in\M_2(\C)$ admits unitaries $U,V$ with
\begin{equation}\label{eq:mi07-target}
 \maxmod(A+B)\leq C\bigl(U\maxmod(A)U^*+V\maxmod(B)V^*\bigr).
\end{equation}
In particular, the proposed constant $C=1$ is false.
\end{theorem}

\begin{proof}
Let $t>0$ and take
\[
 A=e_1e_1^*=\begin{pmatrix}1&0\\0&0\end{pmatrix},\qquad
 B=te_1e_2^*=\begin{pmatrix}0&t\\0&0\end{pmatrix}.
\]
Since $A$ is a projection,
\[
 (|A|^r+|A^*|^r)^{1/r}=2^{1/r}A,
 \qquad\maxmod(A)=A.
\]
For $B$, the right and left moduli are $te_2e_2^*$ and $te_1e_1^*$, so
\[
 \maxmod(B)=tI_2.
\]

Put $s=\sqrt{1+t^2}$ and $v=(e_1+te_2)/s$. The rank-one matrix $A+B=e_1(e_1+te_2)^*$ has moduli
\[
 |A+B|=svv^*,\qquad |(A+B)^*|=se_1e_1^*.
\]
The matrix $R=vv^*+e_1e_1^*$ is positive definite because $v$ and $e_1$ are linearly independent. Consequently
\[
 \bigl(|A+B|^r+|(A+B)^*|^r\bigr)^{1/r}
 =sR^{1/r}\longrightarrow sI_2,
\]
and therefore $\maxmod(A+B)=sI_2$.

For arbitrary $U,V$, the matrix
\[
 U\maxmod(A)U^*+V\maxmod(B)V^*
 =Ue_1e_1^*U^*+tI_2
\]
has eigenvalues $1+t$ and $t$. Taking a unit vector orthogonal to $Ue_1$ in \eqref{eq:mi07-target} forces
\[
 s\leq Ct,\qquad C\geq\frac{\sqrt{1+t^2}}{t}.
\]
This lower bound is unbounded as $t\downarrow0$. For $C=1$, it is already impossible for every $t>0$.
\end{proof}

\begin{remark}
One may choose $t=1/2$ for a fixed explicit counterexample: $\maxmod(A+B)=(\sqrt5/2)I_2$, whereas every two-orbit sum on the right before multiplication by $C$ has eigenvalues $3/2$ and $1/2$. The proof also answers negatively the dimension-dependent finite-constant question preceding the conjecture \cite[Question 3.11]{BL2026}. The calculation uses the spectral-order join, not the usual pointwise maximum of entries or a different symmetric modulus.
\end{remark}
